\begin{abstract}
Retrieval bounds the online token set but recomputes every model layer, whereas
state reuse avoids computation by retaining layer-wise caches. We introduce
\textbf{CoMem}, a cross-query memory interface that stores one intermediate
residual per context token, selects a bounded set of chunks for each query, and
continues computation only through the upper transformer layers. The split
depth exposes an explicit quality--latency--storage trade-off: on a matched
selected pack, resuming at layer 12 yields a $1.403\times$ Read-phase speedup
over raw-text replay while reducing a 15-cell RULER macro by 3.12 points. A
rank-32 self-distillation adapter recovers most of the otherwise severe
interface loss without updating the backbone. For repeated-query workloads, the
one-time Write is amortized after roughly 9 queries at 32k and 26--28 queries at
128k in our serving measurements. The combined bounded-selection operating
point keeps online prefill near 1.1\,s and 18.7\,GB at 128k, versus 71.4\,s and
50.0\,GB for dense prefill on the same H20; this $64.9\times$ figure includes
the benefit of selecting a fixed working set and is not a depth-only effect.
On a paired multikey diagnostic, controlled interface ablations identify
missing lower-layer document context during independent chunk writing as a
major source of fidelity loss. A deployable overlap-Write variant raises that
diagnostic from 92.5 to
98.5--99.0 while leaving persistent storage and Read cost unchanged, at the
price of additional one-time Write compute. CoMem therefore provides a measured
computation-reuse frontier for long-context serving, rather than uniform quality
dominance over raw-text retrieval.
\end{abstract}
